\scanpage{38}
\uline{Autom d'ordre deux d'une sphère holomorphe \inserted{$\Sigma$} à
\inserted{un $K$ de} $5$ pts marqués} : Un pt fixe exactement dans $K$,
soit $K=\{P\}\amalg I$, et $\Sigma$ est munie de $I$ et d'un autom
\inserted{$\sigma$} d'ordre $2$, stabilisant $I$,
\inserted{$I$ ne contenant pas de pt fixe de $\sigma$.}
Donc cel. revient à la donnée

a) D'une sphère hol. $\Sigma$, un autom \inserted{$\sigma$} d'ordre $2$
et un pt fixe marqué --- appelant celui-ci $\infty$, et l'autre pt fixe $0$,
cela revient à la donnée d'un torseur $T_1\ (= \Sigma\setminus\{0,\infty\})$
sous $\mathbb C^*$, i.e. d'une droite vectorielle complexe $L_1$
$(T_1=L_1^*)$

b) D'une paire d'orbites \inserted{(distinctes)} de
$\{\pm1\}\subset\mathbb C^*$ opérant sur $T_1=L_1^*$.
Introduisant $T_2=T^{\wedge2}\simeq L_2^*$, où
$L_2=L_1^{\otimes2}$, cel revient à la donnée d'une paire de pts
\inserted{$\{\xi,\eta\}$} (distincts) de $T_2$.
Considérons alors le pt
\[
 e=\xi\wedge\eta=\eta\wedge\xi
 \in T_4=T^{\wedge4}\simeq L_4^*
 \quad(\text{où }L_4=L_1^{\otimes4})
\]
on a donc un
\[
 \mathrm{iso.}\quad T^{\wedge4}\simeq\mathbf1_{\mathbb C^*}
\]
d'où une réduction du groupe structural \inserted{de $T$}
de $\mathbb C^*$ : $\mu_4(\mathbb C^*)\simeq\mathbb Z/4\mathbb Z$.
On peut dire que la donnée de $\{\xi,\eta\}$ revient à celle de
$e=\xi\wedge\eta\inserted{\in T_4}$ et de
$\xi+\eta\in\unsure{L_2}$, soumis à la condition que
$(\xi+\eta)^2-4\xi\eta\ne0$ \deleted{\unsure{un}}.
Notons \deleted{$\unsure{\xi^2+\eta^2-2\xi\eta\ne}$}
que $T_2$ admet une réduction du groupe structural $\mathbb C^*$ :
$\{\pm1\}$, par la \unsure{bibase} $\pm\sqrt e$, donc $\xi+\eta$
définit un \unsure{élt} de $\mathbb C$ à un signe près, dont
\deleted{$(\xi+\eta)^2$}
\mspage{37}
\unsure{\texttt{00 00}} \illegible{4} \unsure{\texttt{CLJ}}
\illegible{3}

\unsure{\texttt{5 12 15 SIGNES LIGNES}}

\unsure{\texttt{FM}} \illegible{2} \unsure{\texttt{LIGNES}}

\unsure{\texttt{CTRL}}

\unsure{\texttt{DEBUT DU TRAITEMENT}}

\unsure{\texttt{FIN DU TRAITEMENT}}

\unsure{\texttt{DUREE}} \illegible{1} \unsure{\texttt{SEC}}

\unsure{\texttt{2}} \illegible{2}

\unsure{\texttt{0}}

\unsure{\texttt{16.18 TOTL}} \illegible{1}
\unsure{\texttt{880 30.53 45.63 PLUM 0}}

\unsure{\texttt{16.852 4}}

\unsure{\texttt{16.52 1}}

\unsure{\texttt{19.11 TOTL949 1840 74.212 40.33 PLUME 0}}

\unsure{\texttt{19.426 4}}
